% Diagrama conceptual de las tres capas del presente
% Compilar independientemente o \input en el manuscrito
\begin{tikzpicture}[
  node distance=0.6cm,
  layer/.style={rectangle, draw=conceptcolor, 
               minimum width=7.8cm, minimum height=1.1cm, 
               align=center, font=\small, text width=7.5cm,
               inner sep=4pt},
  iarrow/.style={-{Stealth[length=2.5mm]}, thick, conceptcolor!75},
  farrow/.style={-{Stealth[length=2.5mm]}, thick, dashed, conceptcolor!60},
  flabel/.style={font=\footnotesize, text=gray!75!black},
  alabel/.style={font=\footnotesize, text=gray!70!black, align=left}
]
% Title
\node[font=\bfseries\small, text=conceptcolor] at (0, 3.8) {Diagrama conceptual de las tres capas del presente};

% Capa 3 (top)
\node[layer, fill=orange!12!white, draw=orange!70!black] (c3) at (0, 2.9) {
  \textbf{Capa 3: Reorganización estructural del sistema}\\[0.08em]
  \small Norma de Frobenius; picos hiperestructurales
};

% Capa 2
\node[layer, fill=green!8!white, draw=green!55!black] (c2) at (0, 1.5) {
  \textbf{Capa 2: Coherencia entre presentes locales}\\[0.08em]
  \small Reducción de antisincronización $A(k)$
};

% Capa 1 (bottom)
\node[layer, fill=blue!8!white, draw=conceptcolor] (c1) at (0, 0.1) {
  \textbf{Capa 1: Presente local intensificado}\\[0.08em]
  \small Hiper-persistencia + trapping (RQA Config B)
};

% Integration arrows (right side, upward solid)
\draw[iarrow, ->] ([xshift=0.6cm]c1.east) -- ++(0.8,0) |- node[near start, flabel, right, yshift=0.1cm] {requiere $\downarrow A(k)$} ([xshift=0.6cm]c2.east);
\draw[iarrow, ->] ([xshift=0.6cm]c2.east) -- ++(0.8,0) |- node[near start, flabel, right, yshift=0.1cm] {condición necesaria} ([xshift=0.6cm]c3.east);

% Filtering arrows (left side, downward dashed)
\draw[farrow, ->] ([xshift=-0.6cm]c2.west) -- ++(-0.8,0) |- node[near start, flabel, left, yshift=-0.1cm] {filtrado ascendente} ([xshift=-0.6cm]c1.west);
\draw[farrow, ->] ([xshift=-0.6cm]c3.west) -- ++(-0.8,0) |- node[near start, flabel, left, yshift=-0.1cm] {filtrado descendente} ([xshift=-0.6cm]c2.west);

% Side annotation (antisincronización)
\node[alabel, anchor=east] at (-4.8, 1.5) {
  \textbf{Antisincronización $A(k)$}\\[0.05em]
  Alta $\rightarrow$ fragmentación\\[0.05em]
  Baja $\rightarrow$ permeabilidad
};
\end{tikzpicture}